\section{Foundations of the constructible hierarchy}
\label{sec:p06-constructible-hierarchy-foundations}

The following derivation works in the theory \(\mathsf{ZF}\) fixed in
\cref{sec:zf-zfc-01}.  Its inputs are the definitions of transitive set,
ordinal, function, and satisfaction in
\cref{def:transitive-set,def:ordinal-cardinal,sec:zf-zfc-02,sec:fol-semantics-02}.
It supplies the recursion and definability details needed for
\cref{thm:p06-constructible-class}.  Class notation throughout abbreviates
formulas with set parameters.  The arguments below do not use Choice.

\subsection{Ordinals and definable recursion}

\begin{hbdefinition}{def:p06-bounded-ordinal-formula}{A bounded formula for ordinalhood}
A formula of \(\Lang_{\in}\) is \emph{bounded} if its quantifiers have
the forms \(\exists u\in v\) and \(\forall u\in v\), where \(u\ne v\);
these abbreviate the usual guarded quantifiers.  Atomic formulas and
propositional combinations of bounded formulas are bounded.  Write
\[
\begin{aligned}
 \operatorname{Tr}(x)&\quad:\Longleftrightarrow\quad
       \forall u\in x\,\forall v\in u\;(v\in x),\\
 \operatorname{Ord}_{\mathrm b}(x)&\quad:\Longleftrightarrow\quad
       \operatorname{Tr}(x)
       \land\forall u\in x\,\forall v\in u\,\forall w\in v\;(w\in u)\\
       &\hspace{35mm}\land
       \forall u\in x\,\forall v\in x\;
                (u=v\lor u\in v\lor v\in u).
\end{aligned}
\]
Expanding \(\operatorname{Tr}(x)\) gives a bounded formula; bound
variables are renamed freshly when abbreviations are expanded.
The subscript distinguishes this formula from the class
\(\operatorname{Ord}\), defined by the original ordinal definition.
\end{hbdefinition}

\begin{hbtheorem}{thm:p06-ordinal-foundation-facts}{proved=annex:p06.ordinal-foundation-facts}{Ordinal facts needed for hierarchy recursion}
In \(\mathsf{ZF}\), a set is an ordinal exactly when it satisfies
\(\operatorname{Ord}_{\mathrm b}\).  Every member of an ordinal is an
ordinal.  For ordinals \(\alpha,\beta\), exactly one of
\(\alpha\in\beta\), \(\alpha=\beta\), and \(\beta\in\alpha\) holds,
and
\[
 \alpha\subseteq\beta
 \quad\Longleftrightarrow\quad
 \alpha\in\beta\text{ or }\alpha=\beta.
\]
The empty set, every successor \(\alpha+1\), and the union of any set
of ordinals are ordinals.  If \(\lambda\) is a nonzero limit ordinal,
then \(\bigcup\lambda=\lambda\), and \(\beta+1<\lambda\) whenever
\(\beta<\lambda\).  The least inductive set \(\omega\) is an ordinal;
its members other than zero are successors.
\end{hbtheorem}

\begin{hbdefinition}{def:p06-recursion-approximation}{Set approximations to a definable recursion}
Fix a formula \(\vartheta(f,y,\bar p)\) such that \(\mathsf{ZF}\) proves:
for every function \(f\) with ordinal domain there is exactly one
set \(y\) satisfying \(\vartheta(f,y,\bar p)\).
Write \(G(f)=y\) for this formula-defined value.  A
\(G\)-\emph{approximation of length} \(\gamma\) is a function \(f\)
with ordinal domain \(\gamma\) satisfying
\[
 \forall\xi\in\gamma\;
       \vartheta(f\mathbin{\upharpoonright}\xi,f(\xi),\bar p).
\]
Denote this single first-order condition by
\(\operatorname{App}_G(\gamma,f)\).  Function evaluation and restriction
are eliminable set-theoretic abbreviations; the subscript \(G\) records
the fixed formula \(\vartheta\), not a class parameter.
\end{hbdefinition}

\begin{hbtheorem}{thm:p06-definable-ordinal-recursion}{proved=annex:p06.definable-ordinal-recursion}{Definable recursion with set restrictions}
For every ordinal \(\gamma\), there is a unique set function
\(f_\gamma\) satisfying \(\operatorname{App}_G(\gamma,f_\gamma)\).
If \(\beta\leq\gamma\), then
\(f_\beta=f_\gamma\mathbin{\upharpoonright}\beta\).  The formula
\[
 F(\alpha)=y
 \quad:\Longleftrightarrow\quad
 \alpha\in\operatorname{Ord}\ \land
 \exists f\,
   (\operatorname{App}_G(\alpha+1,f)\land f(\alpha)=y)
\]
defines the unique class function satisfying
\(F(\alpha)=G(F\mathbin{\upharpoonright}\alpha)\) on the ordinals.
Its restriction to every set of ordinals is a set function.
\end{hbtheorem}

\subsection{A uniform set construction of definability}

\begin{hbdefinition}{def:p06-finite-formula-instructions}{Finite instruction codes for membership formulas}
Fix variables \(v_i\) for \(i\in\omega\).  An auxiliary formula code
is a nonempty finite list \(c\) of triples of natural numbers.  A
record at position \(k\) has one of the following meanings:
\[
\begin{array}{c|l|l}
\text{record}&\text{formula at this position}&\text{conditions}\\
\hline
(0,i,j)&v_i=v_j&i,j\in\omega\\
(1,i,j)&v_i\in v_j&i,j\in\omega\\
(2,0,0)&\bottom&\\
(3,i,0)&\neg\psi_i&i<k\\
(4,i,j)&\psi_i\land\psi_j&i,j<k\\
(5,i,j)&\psi_i\lor\psi_j&i,j<k\\
(6,i,j)&\psi_i\to\psi_j&i,j<k\\
(7,i,j)&\forall v_i\,\psi_j&j<k,\ i\in\omega\\
(8,i,j)&\exists v_i\,\psi_j&j<k,\ i\in\omega.
\end{array}
\]
Here \(\psi_i\) is the formula at the indicated earlier position.
The output formula is the formula at the final position.  Tuples and
lists use the already fixed set coding of pairs and functions.
Let \(\mathcal C\) be the set of lists satisfying these conditions.
Repeated use of an earlier position merely shares a subformula; it
does not add a logical connective or change satisfaction.

For nonempty \(A\), put \(S_A=A^\omega\), the set of all variable
assignments.  For \(s\in S_A\), write \(s[i\mapsto a]\) for its
update at \(i\).  The notation \(\operatorname{Sat}(A,c,s)\) will
denote the truth value obtained by evaluating the instructions in
\((A,\in)\) under \(s\).  The next theorem constructs this relation
as a set and supplies its uniform first-order definition.
\end{hbdefinition}

\begin{hbtheorem}{thm:p06-set-sized-definability}{proved=annex:p06.set-sized-definability}{Set-sized satisfaction and the definability operation}
For every nonempty set \(A\), there is a uniquely determined
relation \(\operatorname{Sat}_A\subseteq\mathcal C\times S_A\)
satisfying the Tarskian clauses for the instruction codes.  A single
formula of \(\Lang_{\in}\), with \(A,c,s\) as set parameters,
defines \(\operatorname{Sat}(A,c,s)\).  Every ordinary membership
formula has an instruction code with the same satisfaction condition;
every instruction code unfolds to such a formula.

The operation \(A\mapsto\operatorname{Def}(A)\) in
\cref{def:p06-def-empty-base} is consequently a formula-defined
class operation with a unique set value for every set \(A\).
For \(A\ne\varnothing\), an equivalent expression for that value is
\[
 \left\{X\in\mathcal{P}(A):
 \begin{gathered}
  \exists c\in\mathcal C\,\exists s\in S_A\,\forall x\in A\\[-2pt]
  \bigl(x\in X\leftrightarrow
       \operatorname{Sat}(A,c,s[0\mapsto x])\bigr)
 \end{gathered}
 \right\}.
\]
For \(A=\varnothing\), the value is \(\{\varnothing\}\);
no satisfaction relation on an empty structure is used.
\end{hbtheorem}

\begin{hbtheorem}{thm:p06-bounded-transitive-absoluteness}{proved=annex:p06.bounded-transitive-absoluteness}{Bounded absoluteness for a transitive set}
Let \(A\ne\varnothing\) be transitive.  For each fixed bounded
formula \(\varphi\) and assignment of its free variables to members
of \(A\), its satisfaction in \((A,\in)\) is equivalent to the
ambient assertion \(\varphi\) with those same values.  In particular,
for \(x\in A\),
\[
 (A,\in)\entails\operatorname{Ord}_{\mathrm b}(x)
 \quad\Longleftrightarrow\quad x\in\operatorname{Ord}.
\]
The display uses the usual parameter-assignment abbreviation.  Its
ambient right side is a formula of set theory, not satisfaction in
a structure whose domain is the universe.
\end{hbtheorem}

\begin{hblemma}{thm:p06-def-transitive-growth}{proved=annex:p06.def-transitive-growth}{Definability preserves transitivity and includes the previous level}
For every set \(A\),
\(\operatorname{Def}(A)\subseteq\mathcal{P}(A)\) and
\(A\in\operatorname{Def}(A)\).  If \(A\) is transitive, then
\[
 A\subseteq\operatorname{Def}(A),
 \qquad \operatorname{Def}(A)\text{ is transitive}.
\]
These are inclusion statements for an input and its next level.
They do not assert that the operation \(\operatorname{Def}\)
preserves every inclusion between arbitrary transitive inputs.
\end{hblemma}

\subsection{Existence, height, and the constructible class}

\begin{hbtheorem}{thm:p06-hierarchy-recursion-existence}{proved=annex:p06.hierarchy-recursion-existence}{Definable existence of the two hierarchies}
The recursions defining \(L_\alpha\) in
\cref{sec:constructibility-forcing-01} and \(V_\alpha\) in
\cref{sec:cumulative-hierarchy-02} determine unique formula-defined
class functions on the ordinals.  Each value and each restriction
to a set of ordinal indices is a set.  Thus
\[
 x\in L\quad:\Longleftrightarrow\quad
   \exists\alpha\in\operatorname{Ord}\;(x\in L_\alpha)
\]
is a first-order formula after the hierarchy abbreviation is expanded.
\end{hbtheorem}

\begin{hbtheorem}{thm:p06-constructible-level-height}{proved=annex:p06.constructible-level-height}{The exact ordinal height of each constructible level}
For all ordinals \(\alpha\) and \(\beta\):
\begin{enumerate}
 \item \(L_\alpha\) is transitive, and
       \(L_\alpha\subsetneq L_{\alpha+1}\);
 \item if \(\alpha\leq\beta\), then
       \(L_\alpha\subseteq L_\beta\);
 \item \(L_\alpha\cap\operatorname{Ord}=\alpha\), meaning that
       the subset of ordinal members of \(L_\alpha\) is exactly
       the set \(\alpha\);
 \item \(\alpha\in L_{\alpha+1}\), and
       \(L_\alpha\subseteq V_\alpha\).
\end{enumerate}
In the successor step, the subset of \(L_\alpha\) defined by
\(\operatorname{Ord}_{\mathrm b}\) is precisely \(\alpha\).
When \(\alpha=0\), its admission at the next level follows instead
from the explicit empty-base clause.
\end{hbtheorem}

\begin{hbcorollary}{thm:p06-constructible-class-foundation}{proved=annex:p06.constructible-class-foundation}{Transitivity and all ordinals in the constructible class}
In \(\mathsf{ZF}\), the formula-defined class \(L\) is transitive and
contains every ordinal.  It is a proper class.
\end{hbcorollary}

These results establish the hierarchy and its height.
The axiom-by-axiom inner-model proof, canonical well-order,
finite-package condensation argument, and GCH derivation follow in
\cref{sec:p06-constructible-inner-model-completion}.
The separately stated arbitrary-limit condensation theorem retains
its own source standing.
